\subsection{\texorpdfstring{Perfect rings of characteristic \(p\)}{Perfect rings of characteristic p}}

DEFINITION 2. --- A ring A of characteristic \(p \neq 0\) is said to be perfect if it is commutative and the mapping \(a \mapsto a^p\) is bijective.

If the ring A is perfect of characteristic p, the mapping \(a \mapsto a^{p^{f}}\) is an automorphism of the ring A for every integer \(f \geqslant 0\) ; the inverse automorphism is denoted by \(a \mapsto a^{1/p^{f}}\) or \(a \mapsto a^{p^{-f}}\) and the image of a subset S of A under this automorphism is written \(S^{1/p^{f}}\) or \(S^{p^{-f}}\) . It is clear that \((a^{p^{e}})^{p^{f}} = a^{p^{e+f}}\) for all \(a \in A\) and any integers e and f (of whatever sign).

Let A be a commutative ring of characteristic p.~For every integer \(f \geqslant 0\) let us write \(n_{f}\) for the kernel of the endomorphism \(a \mapsto a^{p^{f}}\) of the ring A. Then \((\mathfrak{n}_{f})_{f \geqslant 0}\) is an increasing sequence of ideals of A; since every positive integer is majorized by a power of p, the ideal \(n = \bigcup_{f \geqslant 0} n_{f}\) consists of all the

nilpotent elements of A. In particular, if A is perfect, every nilpotent element of A is zero.

DEFINITION 3. --- Let A be a commutative ring of characteristic \(p \neq 0\) . By a perfect closure of A we understand a pair \((\hat{\mathbf{A}}, u)\) where \(\hat{\mathbf{A}}\) is a perfect ring of characteristic \(p\) and \(u\) is a homomorphism of A into \(\hat{\mathbf{A}}\) satisfying the following universal property:

(PC) If B is a perfect ring of characteristic p and v is a homomorphism of A into B, then there exists a unique homomorphism h of \(\hat{A}\) into B such that \(v = h \circ u\) .

The universal property (PC) implies at once the uniqueness of the perfect closure, in the following sense: if \((\hat{\mathbf{A}}, u)\) and \((\hat{\mathbf{A}}', u')\) are two perfect closures of A, then there exists a unique isomorphism h of \(\hat{A}\) onto \(\hat{A}'\) such that \(u' = h \circ u\) (cf.~E, IV, p.~23). We shall now establish the existence of the perfect closure:

THEOREM 3. --- Let A be a commutative ring of characteristic \(p \neq 0\) . There exists a perfect closure \((\hat{\mathbf{A}}, u)\) of A. Moreover, the kernel of \(u\) is the set of all nilpotent elements of A and for each \(x \in \hat{\mathbf{A}}\) there exists an integer \(n \geqslant 0\) such that \(x^{p^n} \in u(\mathbf{A})\) .

For each integer \(n \geqslant 0\) , put \(\mathbf{A}_n = \mathbf{A}\) ; when \(m \geqslant n\) we define a homomorphism \(\pi_{m,n}\) of \(\mathbf{A}_n\) into \(\mathbf{A}_m\) by \(\pi_{m,n}(a) = a^{p^{m-n}}\) . We thus obtain a direct system of rings \((\mathbf{A}_n, \pi_{m,n})\) (I, p.~120); let \(\hat{\mathbf{A}}\) be the direct limit of this system and \(u_n\) the canonical homomorphism of \(\mathbf{A}_n = \mathbf{A}\) into \(\hat{\mathbf{A}}\) ; we also put \(u = u_0\) . By construction of the direct limit the kernel \(n\) of \(u\) is the union of the kernels of the homomorphisms \(\pi_{n,0}: a \mapsto a^{p^n}\) of A into A, thus it consists of all the nilpotent elements of A. The ring \(\hat{A}\) is commutative of characteristic p by Remark 1 of V, p.~3.

The ring \(\hat{\mathbf{A}}\) is the union of the ascending sequence \((u_n(\mathbf{A}))_{n\geq 0}\) of subrings. We have \(u_{n}(\mathbf{A})^{p^{n}} = \mathbf{A}\) ; hence for each \(x\in \hat{\mathbf{A}}\) there exists an integer \(n\geqslant 0\) such that \(x^{p^n}\in u(\mathbf{A})\) . We also have \(u_{n}(\mathbf{A}) = u_{n + 1}(\mathbf{A})^{p}\) , whence \(\hat{\mathbf{A}}^p = \hat{\mathbf{A}}\) . Let \(x\in \hat{\mathbf{A}}\) be such that \(x^{p} = 0\) ; choose an integer \(n\geqslant 1\) and an element \(a\in \mathbf{A}\) such that \(x = u_{n}(a)\) . Then we have \(u_{n - 1}(a) = u_n(a)^p = 0\) ; by definition of the direct limit there exists an integer \(m\geqslant n\) such that \(\pi_{m - 1,n - 1}(a) = 0\) , that is, \(a^{p^{m - n}} = 0\) . We thus have \(\pi_{m,n}(a) = 0\) , whence \(u_{n}(a) = 0\) , that is, \(x = 0\) . Therefore the ring \(\hat{\mathbf{A}}\) is perfect of characteristic \(p\) .

Let \(v\) be a homomorphism of \(A\) into a perfect ring \(B\) of characteristic \(p\) . For every integer \(n \geqslant 0\) , the mapping \(b \mapsto b^{p^n}\) is an automorphism of \(B\) and so there exists a homomorphism \(v_n\) of \(A_n = A\) into \(B\) characterized by \(v(a) = v_n(a)^{p^n}\) . We then have \(v_m \circ \pi_{m,n} = v_n\) for \(m \geqslant n \geqslant 0\) ; by definition of the direct limit there exists a homomorphism \(h\) of \(\hat{A}\) into \(B\) such that \(v_n = h \circ u_n\) for all \(n \geqslant 0\) ; in particular we have \(v = v_0 = h \circ u_0 = h \circ u\) . Finally let \(h'\) be a homomorphism of \(\hat{A}\) into \(B\) such that \(h' \circ u = v\) . Let \(x \in \hat{A}\) ; as we have seen, there exist an integer \(n \geqslant 0\) and an element \(a \in A\) such that \(x^{p^n} = u(a)\) . Then we have

\[
h (x) ^ {p ^ {n}} = h (u (a)) = v (a) = h ^ {\prime} (u (a)) = h ^ {\prime} (x) ^ {p ^ {n}},
\]

and since B is perfect, we find \(h(x) = h'(x)\) . Thus we have \(h' = h\) , and this completes the proof that \((\hat{\mathbf{A}}, u)\) is a perfect closure of A.

PROPOSITION 3. --- Let B be a perfect ring of characteristic p and A a subring of B. Write \(A^{p^{-\infty}} = \bigcup_{f \geqslant 0} A^{p^{-f}}\) and denote by j the canonical injection of A in

\(\mathbf{A}^{p^{-\infty}}\) . Then \(\mathbf{A}^{p^{-\infty}}\) is the smallest perfect subring of \(\mathbf{B}\) containing \(\mathbf{A}\) and \((\mathbf{A}^{p^{-\infty}}, j)\) is a perfect closure of \(\mathbf{A}\) .

For each integer \(f \in \mathbf{Z}\) denote by \(\pi_f\) the automorphism \(b \mapsto b^{p^f}\) of \(\mathbf{B}\) . The sequence of subrings \(\pi_{-f}(\mathbf{A})\) of \(\mathbf{B}\) (for \(f \geqslant 0\) ) is ascending and its union \(\mathbf{A}^{p^{-\infty}}\) is thus a subring of \(\mathbf{B}\) . We have \(\pi_1(\mathbf{A}^{p^{-\infty}}) = \bigcup_{f \geqslant 0} \pi_{-(f-1)}(\mathbf{A}) = \mathbf{A}^{p^{-\infty}}\) , hence

\(A^{p^{-\infty}}\) is a perfect subring of B. Finally let \(B_0\) be a perfect subring of B containing A; for every integer \(f \geqslant 0\) we have \(\pi_{-f}(A) \subset \pi_{-f}(B_0) = B_0\) , whence \(A^{p^{-\infty}} \subset B_0\) .

If \(v\) is a homomorphism of \(\mathbf{A}\) into a perfect ring \(\mathbf{B}'\) of characteristic \(p\) , then for every integer \(f \geqslant 0\) we can define a homomorphism \(h_f\) of \(\pi_{-f}(\mathbf{A})\) into \(\mathbf{B}'\) by \(h_f(\pi_{-f}(a)) = v(a)^{p^{-f}}\) for all \(a \in \mathbf{A}\) . We see at once that \(h_{f + 1}\) agrees with \(h_f\) on \(\pi_{-f}(\mathbf{A})\) ; thus there exists a homomorphism \(h\) of \(\mathbf{A}^{p^{-\infty}}\) into \(\mathbf{B}'\) which induces \(h_f\) on \(\pi_{-f}(\mathbf{A})\) for all \(f \geqslant 0\) and in particular, \(h\) extends \(h_0 = v\) . If \(h'\) is another extension of v to a homomorphism of \(A^{p^{-\infty}}\) into B, the equality \(h' = h\) may be established as in the proof of Th. 3.

